\section{Vícevrstvé neuronové sítě a hluboké učení. Koncepce, aplikace, technologie}

\subsection*{Základní vymezení}

\textbf{Neuronová síť} je model složený z propojených výpočetních jednotek -- neuronů. Každý neuron počítá
vážený součet vstupů, přičte bias a použije aktivační funkci:
\[
z = w^\top x + b,
\qquad
a = \phi(z).
\]

Vícevrstvá neuronová síť skládá takových transformací více za sebou:
\[
x
\rightarrow
h^{(1)}
\rightarrow
h^{(2)}
\rightarrow
\cdots
\rightarrow
\hat y.
\]

\textbf{Hluboké učení} je část strojového učení založená na neuronových sítích s více skrytými vrstvami.
Hluboké sítě se neučí jen přímý vztah mezi vstupem a výstupem, ale postupně se učí reprezentace dat.
Nižší vrstvy zachycují jednoduché rysy, vyšší vrstvy složitější struktury.

\paragraph{Intuice.}
U obrázku mohou první vrstvy detekovat hrany a jednoduché tvary, další vrstvy části objektů a poslední vrstvy
celé objekty. U textu mohou první vrstvy zachycovat tokeny/slova, další vrstvy syntaktické a významové vztahy.

\subsection*{Mělká vs. hluboká síť}

\begin{itemize}
    \item \textbf{Mělká síť} má typicky jednu skrytou vrstvu.
    \item \textbf{Hluboká síť} má více skrytých vrstev.
    \item Více vrstev umožňuje skládat jednoduché nelineární transformace do složitějších funkcí.
    \item Hluboké sítě jsou výkonné hlavně pro nestrukturovaná data: obraz, text, zvuk, video.
\end{itemize}

Obecný tvar dopředné hluboké sítě:
\[
h^{(0)} = x,
\]
\[
z^{(\ell)} = W^{(\ell)}h^{(\ell-1)} + b^{(\ell)},
\qquad
h^{(\ell)} = \phi^{(\ell)}(z^{(\ell)}),
\qquad
\ell=1,\ldots,L.
\]

Výstup:
\[
\hat y = g(z^{(L)}),
\]
kde $g$ závisí na úloze:
\begin{itemize}
    \item regrese: lineární výstup,
    \item binární klasifikace: sigmoid,
    \item vícetřídní klasifikace: softmax.
\end{itemize}

\subsection*{Proč hluboké učení funguje}

Hluboké sítě jsou úspěšné hlavně díky kombinaci několika faktorů:
\begin{itemize}
    \item dostupnost velkých datasetů,
    \item vyšší výpočetní výkon, hlavně GPU/TPU,
    \item lepší optimalizační algoritmy,
    \item lepší inicializace vah,
    \item regularizace a normalizace,
    \item specializované architektury jako CNN, RNN a transformery,
    \item transfer learning a předtrénované modely.
\end{itemize}

\paragraph{Důležité.}
Hluboké učení není vhodné vždy. U malých tabulkových datasetů mohou být lepší jednodušší modely, například
logistická regrese, rozhodovací strom, random forest nebo gradient boosting. Hluboké sítě se vyplatí hlavně
tam, kde je hodně dat a složitá struktura vstupu.

\subsection*{Vícevrstvý perceptron}

Vícevrstvý perceptron, MLP, je základní dopředná neuronová síť. Každý neuron jedné vrstvy je typicky spojen
se všemi neurony následující vrstvy. Takovým vrstvám se říká plně propojené vrstvy:
\[
\text{Dense layer}.
\]

Pro jednu skrytou vrstvu:
\[
h = \phi(W^{(1)}x+b^{(1)}),
\]
\[
\hat y = g(W^{(2)}h+b^{(2)}).
\]

Pro více skrytých vrstev:
\[
h^{(1)} = \phi(W^{(1)}x+b^{(1)}),
\]
\[
h^{(2)} = \phi(W^{(2)}h^{(1)}+b^{(2)}),
\]
\[
\hat y = g(W^{(3)}h^{(2)}+b^{(3)}).
\]

\subsection*{Počet parametrů v hustých vrstvách}

Počet parametrů mezi dvěma plně propojenými vrstvami:
\[
(n_{\text{in}}+1)n_{\text{out}},
\]
kde $+1$ odpovídá biasu.

\paragraph{Příklad MNIST.}
Obrázek má velikost:
\[
28 \times 28 = 784
\]
vstupních hodnot. Pokud má síť jednu skrytou vrstvu s 256 neurony a výstupní vrstvu s 10 neurony, počet
parametrů je:
\[
784\cdot 256 + 256 + 256\cdot 10 + 10 = 203530.
\]

To ukazuje, že počet parametrů rychle roste. Proto se u obrázků nepoužívají pouze husté vrstvy, ale hlavně
konvoluční sítě.

\subsection*{Aktivační funkce}

Aktivační funkce zavádějí nelinearitu. Bez nich by složení více vrstev bylo opět jen lineární transformací.

\subsubsection*{Sigmoid}

\[
\sigma(z)=\frac{1}{1+e^{-z}}.
\]

Použití:
\begin{itemize}
    \item binární klasifikace,
    \item výstup jako pravděpodobnost.
\end{itemize}

Nevýhoda:
\begin{itemize}
    \item pro velká kladná/záporná $z$ má velmi malý gradient, což může vést k problému mizejícího gradientu.
\end{itemize}

\subsubsection*{Tanh}

\[
\tanh(z)=\frac{e^z-e^{-z}}{e^z+e^{-z}}.
\]

Výstup je v intervalu:
\[
(-1,1).
\]

\subsubsection*{ReLU}

\[
\operatorname{ReLU}(z)=\max(0,z).
\]

Výhody:
\begin{itemize}
    \item jednoduchá,
    \item výpočetně rychlá,
    \item často funguje lépe ve skrytých vrstvách než sigmoid/tanh,
    \item zmírňuje problém mizejícího gradientu.
\end{itemize}

Nevýhoda:
\begin{itemize}
    \item pro záporné vstupy má nulový gradient; neuron může přestat učit.
\end{itemize}

\subsubsection*{Softmax}

Softmax se používá pro vícetřídní klasifikaci:
\[
\hat p_k =
\frac{e^{z_k}}
{\sum_{j=1}^{K} e^{z_j}}.
\]

Platí:
\[
0 \leq \hat p_k \leq 1,
\qquad
\sum_{k=1}^{K}\hat p_k=1.
\]

\subsection*{Ztrátové funkce}

Neuronová síť se učí minimalizací ztrátové funkce:
\[
J(\theta)
=
\frac{1}{n}
\sum_{i=1}^n
L(y_i,\hat y_i),
\]
kde $\theta$ označuje všechny váhy a biasy.

\subsubsection*{Regrese}

Pro regresi se často používá MSE:
\[
L(y,\hat y)=(y-\hat y)^2.
\]

Celkově:
\[
MSE=
\frac{1}{n}
\sum_{i=1}^{n}
(y_i-\hat y_i)^2.
\]

\subsubsection*{Binární klasifikace}

Binary cross-entropy:
\[
L(y,\hat p)
=
-\left[
y\log(\hat p)
+
(1-y)\log(1-\hat p)
\right].
\]

\subsubsection*{Vícetřídní klasifikace}

Categorical cross-entropy:
\[
L(y,\hat p)
=
-\sum_{k=1}^{K} y_k\log(\hat p_k),
\]
kde $y_k$ je one-hot reprezentace skutečné třídy.

\subsection*{Trénování hlubokých sítí}

Trénování znamená hledání parametrů, které minimalizují ztrátovou funkci:
\[
\min_{\theta} J(\theta).
\]

Nejčastěji se používají gradientní metody:
\[
\theta^{(t+1)}
=
\theta^{(t)}
-
\eta \nabla J(\theta^{(t)}),
\]
kde $\eta$ je learning rate.

\subsubsection*{Forward propagation}

Dopředný průchod spočítá predikci:
\[
x \rightarrow \hat y.
\]

Například:
\[
z^{(1)}=W^{(1)}x+b^{(1)},
\]
\[
h^{(1)}=\phi(z^{(1)}),
\]
\[
z^{(2)}=W^{(2)}h^{(1)}+b^{(2)},
\]
\[
\hat y=g(z^{(2)}).
\]

\subsubsection*{Backpropagation}

Backpropagation efektivně počítá derivace ztráty podle všech vah pomocí řetězového pravidla:
\[
\frac{\partial L}{\partial w}
=
\frac{\partial L}{\partial \hat y}
\cdot
\frac{\partial \hat y}{\partial z}
\cdot
\frac{\partial z}{\partial w}.
\]

\paragraph{Intuice.}
Síť nejprve udělá predikci, poté spočítá chybu a následně šíří informaci o chybě zpět přes jednotlivé vrstvy.
Každá váha se upraví podle toho, jak přispěla k chybě.

\subsection*{Gradientní optimalizátory}

\subsubsection*{Batch gradient descent}

Používá celý trénovací dataset pro jeden update:
\[
\nabla J(\theta)
=
\frac{1}{n}
\sum_{i=1}^{n}
\nabla L_i(\theta).
\]

Výhoda:
\begin{itemize}
    \item stabilní gradient.
\end{itemize}

Nevýhoda:
\begin{itemize}
    \item pomalé u velkých datasetů.
\end{itemize}

\subsubsection*{Stochastic gradient descent}

Aktualizuje parametry po jednom pozorování:
\[
\theta
\leftarrow
\theta
-
\eta \nabla L_i(\theta).
\]

Výhody:
\begin{itemize}
    \item rychlé aktualizace,
    \item šum v gradientu může pomoci uniknout z horších lokálních minim.
\end{itemize}

Nevýhody:
\begin{itemize}
    \item nestabilnější konvergence,
    \item citlivost na learning rate.
\end{itemize}

\subsubsection*{Mini-batch gradient descent}

Používá malou dávku pozorování:
\[
B\subset \{1,\ldots,n\}.
\]

Aktualizace:
\[
\theta
\leftarrow
\theta
-
\eta
\frac{1}{|B|}
\sum_{i\in B}
\nabla L_i(\theta).
\]

Toto je nejběžnější varianta pro hluboké učení.

\subsubsection*{Momentum}

Momentum používá informaci o předchozích gradientech:
\[
v_t = \gamma v_{t-1} + \eta \nabla J(\theta_t),
\]
\[
\theta_{t+1} = \theta_t - v_t.
\]

Pomáhá zrychlit učení ve stabilním směru a tlumit oscilace.

\subsubsection*{Adam}

Adam kombinuje momentum a adaptivní learning rate. Udržuje odhad prvního a druhého momentu gradientů.
V praxi je jedním z nejpoužívanějších optimalizátorů pro neuronové sítě.

\subsection*{Problémy při trénování}

\subsubsection*{Mizející gradient}

Mizející gradient znamená:
\[
\left|
\frac{\partial L}{\partial w}
\right|
\approx 0.
\]

Potom se váhy v nižších vrstvách téměř neučí. Problém je častý u hlubokých sítí a aktivačních funkcí
sigmoid/tanh.

Řešení:
\begin{itemize}
    \item ReLU a její varianty,
    \item vhodná inicializace vah,
    \item batch normalization,
    \item reziduální spojení,
    \item LSTM/GRU u sekvenčních dat.
\end{itemize}

\subsubsection*{Explodující gradient}

Explodující gradient znamená, že gradienty jsou příliš velké:
\[
\left|
\frac{\partial L}{\partial w}
\right|
\gg 1.
\]

Řešení:
\begin{itemize}
    \item gradient clipping,
    \item menší learning rate,
    \item lepší inicializace,
    \item normalizace.
\end{itemize}

\subsubsection*{Overfitting}

Hluboké sítě mají mnoho parametrů, a proto se mohou přeučit:
\[
Err_{\text{train}} \ll Err_{\text{test}}.
\]

Řešení:
\begin{itemize}
    \item větší dataset,
    \item regularizace,
    \item dropout,
    \item early stopping,
    \item data augmentation,
    \item jednodušší architektura,
    \item transfer learning.
\end{itemize}

\subsection*{Regularizace}

\subsubsection*{L2 regularizace}

\[
J_{\lambda}(\theta)
=
J(\theta)
+
\lambda
\sum_j w_j^2.
\]

L2 penalizuje velké váhy a tím omezuje složitost modelu.

\subsubsection*{L1 regularizace}

\[
J_{\lambda}(\theta)
=
J(\theta)
+
\lambda
\sum_j |w_j|.
\]

L1 může vést k řídkým modelům, protože některé váhy tlačí k nule.

\subsubsection*{Dropout}

Dropout náhodně vypíná část neuronů během trénování:
\[
\tilde h = m \odot h,
\qquad
m_j \sim \operatorname{Bernoulli}(q).
\]

Model se tím učí nebýt příliš závislý na konkrétních neuronech.

\subsubsection*{Batch normalization}

Batch normalization normalizuje aktivace v dávce:
\[
\hat x =
\frac{x-\mu_B}
{\sqrt{\sigma_B^2+\varepsilon}},
\]
\[
y=\gamma \hat x+\beta.
\]

Pomáhá stabilizovat a zrychlit trénování.

\subsubsection*{Early stopping}

Sledujeme validační chybu:
\[
Err_{\text{val}}.
\]

Pokud trénovací chyba klesá, ale validační chyba začne růst, model se začíná přeučovat:
\[
Err_{\text{train}}\downarrow
\quad
\text{a}
\quad
Err_{\text{val}}\uparrow.
\]

Trénování zastavíme v bodě s nejlepší validační chybou.

\subsection*{Konvoluční neuronové sítě}

Konvoluční neuronové sítě, CNN, jsou specializované hlavně na obrazová data. Obrázek lze chápat jako tensor:
\begin{itemize}
    \item černobílý obrázek: výška $\times$ šířka,
    \item barevný RGB obrázek: výška $\times$ šířka $\times$ 3,
    \item dávka obrázků: počet obrázků $\times$ výška $\times$ šířka $\times$ počet kanálů.
\end{itemize}

Například RGB batch:
\[
(256,128,128,3).
\]

\subsection*{Proč nestačí plně propojené vrstvy pro obrázky}

U obrázku $32\times 32\times 3$ máme:
\[
32\cdot 32\cdot 3=3072
\]
vstupů. Pokud by byl každý neuron ve skryté vrstvě propojen se všemi pixely, počet parametrů by rychle
rostl. Navíc by síť nevyužívala prostorovou strukturu obrázku.

CNN používají:
\begin{itemize}
    \item lokální propojení,
    \item sdílení parametrů,
    \item konvoluční filtry,
    \item pooling,
    \item hierarchickou extrakci příznaků.
\end{itemize}

\subsection*{Konvoluce}

Konvoluce aplikuje filtr, také kernel, na vstupní obrázek a vytváří feature mapu:
\[
S(i,j)=(I*K)(i,j).
\]

Pro více kanálů a více filtrů:
\[
S_{i,j,k}
=
\phi
\left(
b_k+
\sum_{u=1}^{F_h}
\sum_{v=1}^{F_w}
\sum_{c=1}^{C_{\text{in}}}
K_{u,v,c,k}X_{i+u,j+v,c}
\right).
\]

Kde:
\begin{itemize}
    \item $X$ je vstupní obrázek nebo aktivace předchozí vrstvy,
    \item $K$ je konvoluční filtr,
    \item $S$ je výstupní feature mapa,
    \item $C_{\text{in}}$ je počet vstupních kanálů,
    \item $C_{\text{out}}$ je počet filtrů, tedy počet výstupních kanálů.
\end{itemize}

\subsection*{Velikost výstupu konvoluce}

Pro výšku:
\[
H_{\text{out}}
=
\left\lfloor
\frac{H+2P-F_h}{s_h}
\right\rfloor
+1.
\]

Pro šířku:
\[
W_{\text{out}}
=
\left\lfloor
\frac{W+2P-F_w}{s_w}
\right\rfloor
+1.
\]

Kde:
\begin{itemize}
    \item $H,W$ jsou rozměry vstupu,
    \item $F_h,F_w$ rozměry filtru,
    \item $P$ padding,
    \item $s_h,s_w$ stride.
\end{itemize}

\subsection*{Počet parametrů v konvoluční vrstvě}

Počet trénovatelných parametrů:
\[
F_hF_wC_{\text{in}}C_{\text{out}}+C_{\text{out}}.
\]

První část jsou váhy filtrů, druhá část biasy.

\paragraph{Příklad.}
Konvoluce s filtrem $3\times 3$, jedním vstupním kanálem a 32 filtry má:
\[
3\cdot 3\cdot 1\cdot 32 + 32 = 320
\]
parametrů.

\subsection*{Pooling}

Pooling zmenšuje prostorové rozměry:
\[
H\times W \rightarrow H'\times W'.
\]

Nejčastější je max pooling, například filtr $2\times 2$ se stride 2:
\[
\max
\begin{pmatrix}
x_{11} & x_{12}\\
x_{21} & x_{22}
\end{pmatrix}.
\]

Pooling:
\begin{itemize}
    \item snižuje počet hodnot,
    \item zmenšuje výpočetní náročnost,
    \item zvyšuje určitou robustnost vůči malým posunům,
    \item obvykle nemá trénovatelné parametry.
\end{itemize}

\subsection*{Typická CNN architektura}

Základní konvoluční síť:
\[
\text{Input}
\rightarrow
\text{Convolution}
\rightarrow
\text{ReLU}
\rightarrow
\text{Pooling}
\rightarrow
\text{Convolution}
\rightarrow
\text{ReLU}
\rightarrow
\text{Pooling}
\rightarrow
\text{Flatten}
\rightarrow
\text{Dense}
\rightarrow
\text{Softmax}.
\]

\paragraph{Příklad MNIST CNN.}
\[
(28,28,1)
\rightarrow
(26,26,32)
\rightarrow
(13,13,32)
\rightarrow
(11,11,64)
\rightarrow
(5,5,64)
\rightarrow
(3,3,128)
\rightarrow
(1152,)
\rightarrow
(10,).
\]

Počet parametrů:
\[
3\cdot3\cdot1\cdot32+32
+
3\cdot3\cdot32\cdot64+64
+
3\cdot3\cdot64\cdot128+128
+
3\cdot3\cdot128\cdot10+10
=
104202.
\]

\subsection*{Data augmentation}

Data augmentation uměle rozšiřuje trénovací data transformacemi, které nemění třídu objektu.

U obrázků:
\begin{itemize}
    \item otočení,
    \item posun,
    \item změna měřítka,
    \item ořez,
    \item horizontální překlopení,
    \item šum,
    \item změna jasu nebo kontrastu.
\end{itemize}

\paragraph{Intuice.}
Pes zůstává psem i po mírném otočení nebo ořezu obrázku. Model se tím učí robustnější reprezentace.

Existují dva přístupy:
\begin{itemize}
    \item \textbf{offline augmentation} -- nové obrázky se předem uloží na disk,
    \item \textbf{on-the-fly augmentation} -- transformace se generují během trénování.
\end{itemize}

\subsection*{Moderní CNN architektury}

\begin{itemize}
    \item \textbf{LeNet} -- raná CNN pro rozpoznávání číslic.
    \item \textbf{AlexNet} -- významný úspěch v počítačovém vidění, využití hlubší sítě a GPU.
    \item \textbf{VGG} -- jednoduchá architektura s mnoha $3\times 3$ konvolucemi.
    \item \textbf{Inception / GoogLeNet} -- paralelní konvoluce různých velikostí, efektivní využití parametrů.
    \item \textbf{ResNet} -- reziduální spojení, umožňuje trénovat velmi hluboké sítě.
\end{itemize}

Reziduální blok:
\[
h^{(\ell+1)}
=
F(h^{(\ell)}) + h^{(\ell)}.
\]

Reziduální spojení pomáhá toku gradientu a řeší problém, že velmi hluboké sítě se jinak obtížně trénují.

\subsection*{Rekurentní neuronové sítě}

Rekurentní neuronové sítě, RNN, jsou určeny pro sekvenční data:
\[
x_1,x_2,\ldots,x_T.
\]

Příklady:
\begin{itemize}
    \item text,
    \item řeč,
    \item časové řady,
    \item DNA sekvence,
    \item video.
\end{itemize}

U sekvenčních dat záleží na pořadí. RNN umožňuje, aby zpracování aktuálního prvku záviselo na předchozí
historii.

\subsection*{Vstupní tensor u RNN}

Pro jednu sekvenci:
\[
(x_1,\ldots,x_T),
\]
kde $x_t$ je vektor příznaků v čase $t$.

Pro mini-batch:
\[
(\text{batch size},\text{time steps},\text{features}).
\]

\subsection*{Jednoduchá RNN}

Skrytý stav:
\[
h_t = f_h(W_hh_{t-1}+W_xx_t+b).
\]

Výstup:
\[
y_t = f_o(W_oh_t+c).
\]

Kde:
\begin{itemize}
    \item $h_t$ shrnuje dosavadní historii,
    \item $x_t$ je aktuální vstup,
    \item stejné váhy se používají ve všech časových krocích.
\end{itemize}

\paragraph{Sdílení parametrů.}
RNN používá stejné matice $W_h,W_x,W_o$ pro každý časový krok. Díky tomu může zpracovávat sekvence různé délky.

\subsection*{Typy RNN úloh}

\begin{itemize}
    \item \textbf{One-to-one:}
    jeden vstup a jeden výstup, běžná dopředná síť.
    
    \item \textbf{Many-to-one:}
    sekvence na jeden výstup, například sentiment věty.
    
    \item \textbf{One-to-many:}
    jeden vstup na sekvenci výstupů, například generování popisu obrázku.
    
    \item \textbf{Many-to-many:}
    sekvence na sekvenci, například strojový překlad nebo značkování slov ve větě.
\end{itemize}

\subsection*{Problémy jednoduchých RNN}

Jednoduché RNN mají problém s dlouhodobými závislostmi.

\begin{itemize}
    \item \textbf{Vanishing gradient:}
    informace z dávné minulosti se při zpětném šíření ztrácí.
    
    \item \textbf{Exploding gradient:}
    gradienty mohou prudce růst.
    
    \item \textbf{Sekvenční výpočet:}
    časové kroky se zpracovávají postupně, což omezuje paralelizaci.
\end{itemize}

Řešení:
\begin{itemize}
    \item LSTM,
    \item GRU,
    \item bidirectional RNN,
    \item attention mechanismus,
    \item transformery.
\end{itemize}

\subsection*{LSTM}

LSTM, Long Short-Term Memory, používá paměťovou buňku a brány, které řídí tok informací.

Forget gate:
\[
f_t = \sigma(W_f[h_{t-1},x_t]+b_f).
\]

Input gate:
\[
i_t = \sigma(W_i[h_{t-1},x_t]+b_i).
\]

Kandidátní paměť:
\[
\tilde c_t = \tanh(W_c[h_{t-1},x_t]+b_c).
\]

Aktualizace paměti:
\[
c_t = f_t\odot c_{t-1}+i_t\odot \tilde c_t.
\]

Output gate:
\[
o_t = \sigma(W_o[h_{t-1},x_t]+b_o).
\]

Nový skrytý stav:
\[
h_t = o_t\odot \tanh(c_t).
\]

\paragraph{Intuice.}
LSTM se učí, co zapomenout, co přidat do paměti a co poslat dál jako výstup. Díky tomu lépe zachycuje
dlouhodobé závislosti než jednoduchá RNN.

\subsection*{GRU}

GRU, Gated Recurrent Unit, je jednodušší varianta bránové RNN.

Update gate:
\[
z_t = \sigma(W_z[h_{t-1},x_t]+b_z).
\]

Reset gate:
\[
r_t = \sigma(W_r[h_{t-1},x_t]+b_r).
\]

Kandidátní stav:
\[
\tilde h_t =
\tanh(W_h[r_t\odot h_{t-1},x_t]+b_h).
\]

Nový stav:
\[
h_t =
(1-z_t)\odot h_{t-1}
+
z_t\odot \tilde h_t.
\]

GRU má méně parametrů než LSTM, a proto může být rychlejší.

\subsection*{Bidirectional RNN}

Bidirectional RNN zpracovává sekvenci oběma směry:
\[
\overrightarrow{h_t}
\quad \text{z minulosti do budoucnosti},
\]
\[
\overleftarrow{h_t}
\quad \text{z budoucnosti do minulosti}.
\]

Výstup může používat:
\[
h_t =
[\overrightarrow{h_t},\overleftarrow{h_t}].
\]

Použití:
\begin{itemize}
    \item rozpoznávání řeči,
    \item rozpoznávání rukopisu,
    \item bioinformatika,
    \item zpracování textu, kde známe celou větu.
\end{itemize}

\subsection*{Word embeddings}

Text nelze přímo vložit do neuronové sítě jako slova. Je potřeba převod slov na vektory.

\subsubsection*{One-hot encoding}

Každé slovo je reprezentováno dlouhým řídkým vektorem:
\[
\text{pes} = (0,0,1,0,\ldots,0).
\]

Nevýhody:
\begin{itemize}
    \item vysoká dimenze,
    \item řídkost,
    \item žádná informace o významové podobnosti slov.
\end{itemize}

\subsubsection*{Embedding}

Embedding je hustý nízkorozměrný vektor:
\[
\text{pes} \mapsto e_{\text{pes}}\in \mathbb{R}^d.
\]

Podobná slova by měla mít podobné vektory. Podobnost lze měřit například kosinovou podobností:
\[
\cos(e_i,e_j)=
\frac{e_i^\top e_j}
{\|e_i\|\|e_j\|}.
\]

\subsubsection*{Word2Vec}

Word2Vec se učí reprezentace slov z kontextu.

\begin{itemize}
    \item \textbf{CBOW:}
    z okolních slov predikuje prostřední slovo.
    
    \item \textbf{Skip-gram:}
    z daného slova predikuje okolní slova.
\end{itemize}

Skip-gram bývá vhodnější pro méně častá slova a menší trénovací data. CBOW se obvykle trénuje rychleji.

\subsection*{Transformery}

Transformery jsou architektury založené na mechanismu attention. Vznikly jako řešení omezení RNN:
\begin{itemize}
    \item RNN zpracovávají sekvenci postupně,
    \item obtížně zachycují dlouhé závislosti,
    \item trpí mizejícím gradientem,
    \item hůře se paralelizují.
\end{itemize}

Transformer reprezentuje prvky sekvence podle jejich vzájemné relevance.

\subsection*{Self-attention}

Máme sekvenci reprezentovanou maticí:
\[
X\in \mathbb{R}^{n\times d}.
\]

Z ní se vytvoří:
\[
Q=XW_Q,
\qquad
K=XW_K,
\qquad
V=XW_V.
\]

Kde:
\begin{itemize}
    \item $Q$ jsou queries,
    \item $K$ jsou keys,
    \item $V$ jsou values.
\end{itemize}

Scaled dot-product attention:
\[
\operatorname{Attention}(Q,K,V)
=
\operatorname{softmax}
\left(
\frac{QK^\top}{\sqrt{d_k}}
\right)
V.
\]

\paragraph{Intuice.}
Každý token se ptá, na které ostatní tokeny se má dívat. Podobnost query a key určuje váhu. Výsledkem je
vážený součet value vektorů.

\subsection*{Multi-head attention}

Místo jedné attention se použije více hlav:
\[
\operatorname{head}_i
=
\operatorname{Attention}
(QW_i^Q,KW_i^K,VW_i^V).
\]

Výstup:
\[
\operatorname{MultiHead}(Q,K,V)
=
\operatorname{Concat}(\operatorname{head}_1,\ldots,\operatorname{head}_h)W^O.
\]

Každá hlava se může naučit jiný typ vztahu v sekvenci.

\subsection*{Pozicové kódování}

Self-attention sama o sobě nezná pořadí tokenů. Proto se přidává pozicové kódování:
\[
X_{\text{input}} = X_{\text{embedding}} + X_{\text{position}}.
\]

Bez pozicového kódování by transformer nerozlišil, zda slovo stojí na začátku nebo na konci věty.

\subsection*{Encoder-decoder transformer}

Původní transformer pro strojový překlad má dvě části:

\begin{itemize}
    \item \textbf{Encoder:}
    převádí vstupní sekvenci na kontextové reprezentace.
    
    \item \textbf{Decoder:}
    generuje výstupní sekvenci autoregresivně, tedy postupně token po tokenu.
\end{itemize}

V encoderu se používá self-attention nad vstupem. V decoderu se používá:
\begin{itemize}
    \item masked self-attention, aby model neviděl budoucí tokeny,
    \item attention nad výstupem encoderu,
    \item feed-forward vrstvy.
\end{itemize}

\subsection*{Typy transformerových modelů}

\begin{itemize}
    \item \textbf{Encoder-only:}
    vhodné pro porozumění textu, klasifikaci, extrakci informací.
    Příklad: BERT.
    
    \item \textbf{Decoder-only:}
    vhodné pro generování textu.
    Příklad: GPT modely.
    
    \item \textbf{Encoder-decoder:}
    vhodné pro překlad, sumarizaci a transformaci jedné sekvence na druhou.
    Příklad: T5.
\end{itemize}

\subsection*{Předtrénování a transfer learning}

Hluboké modely často netrénujeme od nuly. Používáme předtrénované modely:
\[
\theta_{\text{pretrained}}.
\]

Poté je upravíme pro konkrétní úlohu:
\[
\theta_{\text{fine-tuned}}.
\]

\subsubsection*{Transfer learning u CNN}

U obrázků se často použije předtrénovaná CNN jako extraktor příznaků:
\[
x \rightarrow \text{CNN backbone} \rightarrow z.
\]

Potom přidáme novou klasifikační hlavu:
\[
z \rightarrow \text{Dense} \rightarrow \hat y.
\]

Možnosti:
\begin{itemize}
    \item zmrazit původní vrstvy a trénovat jen novou hlavu,
    \item částečně odemrazit vyšší vrstvy,
    \item fine-tunovat celý model s malým learning rate.
\end{itemize}

\subsubsection*{Self-supervised learning}

U velkých jazykových modelů se často využívá self-supervised učení. Model si vytváří cílovou proměnnou z
dat samotných.

Příklady:
\begin{itemize}
    \item predikce dalšího tokenu,
    \item predikce maskovaného tokenu,
    \item kontrastivní učení reprezentací.
\end{itemize}

\subsection*{Autoencodery}

Autoencoder je neuronová síť, která se učí rekonstruovat vlastní vstup:
\[
x \rightarrow z \rightarrow \hat x.
\]

Encoder:
\[
z=f_{\theta}(x).
\]

Decoder:
\[
\hat x=g_{\phi}(z).
\]

Ztráta:
\[
L(x,\hat x)=\|x-\hat x\|^2.
\]

\paragraph{Bottleneck.}
Nejužší skrytá vrstva se nazývá code nebo bottleneck. Pokud má malou dimenzi, autoencoder se musí naučit
kompaktní reprezentaci dat.

Použití:
\begin{itemize}
    \item redukce dimenze,
    \item vizualizace dat,
    \item denoising,
    \item detekce anomálií,
    \item inicializace reprezentací.
\end{itemize}

\subsection*{Denoising autoencoder}

Denoising autoencoder dostane zašuměný vstup:
\[
\tilde x = x + \varepsilon,
\]
ale učí se rekonstruovat původní vstup:
\[
\hat x \approx x.
\]

Tím se učí robustnější reprezentace.

\subsection*{Aplikace hlubokého učení}

\subsubsection*{Obrazová data}

\begin{itemize}
    \item klasifikace obrázků,
    \item detekce objektů,
    \item segmentace obrazu,
    \item OCR,
    \item medicínské snímky,
    \item detekce porušování lidských práv z obrazových dat,
    \item kontrola kvality ve výrobě,
    \item autonomní řízení.
\end{itemize}

\subsubsection*{Text a jazyk}

\begin{itemize}
    \item sentiment analysis,
    \item strojový překlad,
    \item sumarizace,
    \item vyhledávání informací,
    \item otázky a odpovědi,
    \item klasifikace dokumentů,
    \item extrakce entit,
    \item generování textu.
\end{itemize}

\subsubsection*{Zvuk a řeč}

\begin{itemize}
    \item rozpoznávání řeči,
    \item převod textu na řeč,
    \item detekce řečníka,
    \item klasifikace zvuků,
    \item zpracování hudby.
\end{itemize}

\subsubsection*{Časové řady}

\begin{itemize}
    \item predikce poptávky,
    \item predikce spotřeby energie,
    \item detekce anomálií,
    \item prediktivní údržba,
    \item finanční časové řady,
    \item senzory a IoT.
\end{itemize}

\subsubsection*{Business aplikace}

\begin{itemize}
    \item detekce fraudu,
    \item scoring zákazníků,
    \item churn prediction,
    \item doporučovací systémy,
    \item automatizace zákaznické podpory,
    \item personalizovaný marketing,
    \item analýza clickstreamů.
\end{itemize}

\subsection*{Technologie pro hluboké učení}

\subsubsection*{Programovací jazyky a knihovny}

Nejčastější technologie:
\begin{itemize}
    \item \textbf{Python} -- hlavní jazyk pro deep learning,
    \item \textbf{NumPy, pandas} -- práce s daty,
    \item \textbf{scikit-learn} -- klasické ML a preprocessing,
    \item \textbf{TensorFlow / Keras} -- tvorba a trénování neuronových sítí,
    \item \textbf{PyTorch} -- flexibilní framework, často používaný ve výzkumu i produkci,
    \item \textbf{Hugging Face Transformers} -- předtrénované NLP a transformerové modely,
    \item \textbf{OpenCV} -- počítačové vidění.
\end{itemize}

\subsubsection*{Hardware}

Hluboké sítě jsou výpočetně náročné. Používají se:
\begin{itemize}
    \item CPU pro menší úlohy,
    \item GPU pro paralelní výpočty matic a tensorů,
    \item TPU pro specializované trénování neuronových sítí,
    \item distribuované clustery pro velké modely.
\end{itemize}

GPU zrychluje zejména operace:
\[
XW,
\qquad
\text{konvoluce},
\qquad
\text{backpropagation}.
\]

\subsubsection*{Datové pipeline}

Pro úspěšné hluboké učení nestačí jen model. Je potřeba i spolehlivá datová pipeline:
\begin{itemize}
    \item načítání dat,
    \item čištění dat,
    \item transformace,
    \item augmentace,
    \item batching,
    \item caching,
    \item monitoring kvality dat.
\end{itemize}

\subsubsection*{Nasazení modelu}

Model lze nasadit například jako:
\begin{itemize}
    \item REST API,
    \item batch scoring,
    \item stream processing,
    \item model běžící v mobilní aplikaci,
    \item model v cloudové službě.
\end{itemize}

Produkční řešení typicky vyžaduje:
\begin{itemize}
    \item kontejnerizaci,
    \item modularizaci,
    \item škálování,
    \item správu přístupů a tajemství,
    \item CI/CD,
    \item auditovatelnost,
    \item možnost rollbacku,
    \item monitoring výkonu a driftu.
\end{itemize}

\subsection*{MLOps}

MLOps propojuje vývoj modelu, datové inženýrství a provoz.

Typický proces:
\begin{enumerate}
    \item příprava dat,
    \item trénování modelu,
    \item validace,
    \item registrace modelu,
    \item nasazení,
    \item monitoring,
    \item retrénování,
    \item správa verzí.
\end{enumerate}

\paragraph{Důležité.}
Model v produkci není jednorázový výstup. Časem se mohou měnit data, chování uživatelů i business procesy.
Proto je nutné sledovat:
\begin{itemize}
    \item predikční výkon,
    \item datový drift,
    \item konceptový drift,
    \item latenci,
    \item náklady,
    \item chybovost služby.
\end{itemize}

\subsection*{Evaluace hlubokých modelů}

Model se hodnotí na oddělených datech:
\[
D =
D_{\text{train}}
\cup
D_{\text{validation}}
\cup
D_{\text{test}}.
\]

\begin{itemize}
    \item \textbf{Trénovací data} slouží k učení vah.
    \item \textbf{Validační data} slouží k výběru architektury a hyperparametrů.
    \item \textbf{Testovací data} slouží k finálnímu nezávislému ověření.
\end{itemize}

U klasifikace:
\[
Accuracy,\quad Precision,\quad Recall,\quad F_1,\quad AUC.
\]

U regrese:
\[
MSE,\quad RMSE,\quad MAE.
\]

U generativních modelů je hodnocení obtížnější, protože kromě přesnosti řešíme také kvalitu, užitečnost,
bezpečnost a věcnou správnost výstupu.

\subsection*{Hyperparametry}

Hyperparametry nejsou přímo naučené z dat, ale nastavují se před trénováním.

Typické hyperparametry:
\begin{itemize}
    \item počet vrstev,
    \item počet neuronů,
    \item typ aktivační funkce,
    \item learning rate,
    \item batch size,
    \item počet epoch,
    \item dropout rate,
    \item regularizační koeficient,
    \item typ optimalizátoru,
    \item velikost kernelu u CNN,
    \item počet attention heads u transformeru.
\end{itemize}

Ladění:
\begin{itemize}
    \item grid search,
    \item random search,
    \item Bayesian optimization,
    \item ruční ladění podle learning curves,
    \item early stopping.
\end{itemize}

\subsection*{Interpretovatelnost a vysvětlitelnost}

Neuronové sítě a hluboké modely jsou často \textbf{black-box} modely. Uživatel obvykle nevidí jednoduché
pravidlo, proč model rozhodl právě takto.

Rozlišujeme:
\begin{itemize}
    \item \textbf{globální interpretovatelnost} -- jak model funguje jako celek,
    \item \textbf{lokální interpretovatelnost} -- proč model rozhodl u konkrétního pozorování.
\end{itemize}

Metody vysvětlování:
\begin{itemize}
    \item variable importance,
    \item permutation importance,
    \item partial dependence plots,
    \item ICE plots,
    \item LIME,
    \item SHAP,
    \item surrogate modely,
    \item saliency maps u obrazových modelů,
    \item attention visualization u transformerů.
\end{itemize}

\paragraph{Surrogate model.}
Naučíme jednodušší model, například rozhodovací strom, který aproximuje predikce složitého modelu:
\[
\hat y_{\text{black-box}}
\approx
\hat y_{\text{surrogate}}.
\]

Potom interpretujeme jednodušší model. Je ale nutné kontrolovat fidelitu, tedy jak dobře surrogate skutečně
napodobuje původní model.

\subsection*{Rizika hlubokého učení}

\begin{itemize}
    \item potřeba velkého množství dat,
    \item vysoké výpočetní náklady,
    \item přeučení,
    \item nízká interpretovatelnost,
    \item citlivost na bias v datech,
    \item data leakage,
    \item adversariální příklady,
    \item drift v produkci,
    \item právní a etické požadavky,
    \item riziko nadměrného spoléhání na výstupy AI.
\end{itemize}

U vysoce rizikových aplikací je nutná lidská kontrola, dokumentace, auditovatelnost a možnost zásahu člověka.

\subsection*{Srovnání hlavních architektur}

\begin{itemize}
    \item \textbf{MLP:}
    obecná dopředná síť pro tabulková data a jednoduché úlohy.
    
    \item \textbf{CNN:}
    nejlepší pro obrazová data a úlohy, kde je důležitá lokální struktura.
    
    \item \textbf{RNN/LSTM/GRU:}
    sekvenční modely, vhodné pro text, řeč a časové řady, ale hůře paralelizovatelné.
    
    \item \textbf{Transformer:}
    attention-based architektura, výborná pro jazyk, text, multimodální modely a dlouhé závislosti.
    
    \item \textbf{Autoencoder:}
    model pro učení reprezentací, redukci dimenze, rekonstrukci a detekci anomálií.
\end{itemize}

\subsection*{Co říct u státnic}

Vícevrstvé neuronové sítě skládají více nelineárních transformací za sebou:
\[
h^{(\ell)}
=
\phi(W^{(\ell)}h^{(\ell-1)}+b^{(\ell)}).
\]
Díky tomu se mohou učit složité vztahy a hierarchické reprezentace dat. Hluboké učení je úspěšné hlavně díky
velkým datům, GPU, lepším optimalizačním algoritmům, regularizaci a specializovaným architekturám.

U obrázků se používají CNN, které využívají konvoluce, lokální propojení, sdílení parametrů a pooling.
U sekvenčních dat se používají RNN, LSTM a GRU:
\[
h_t=f_h(W_hh_{t-1}+W_xx_t+b).
\]
Moderní jazykové a generativní modely stojí hlavně na transformerech a self-attention:
\[
\operatorname{Attention}(Q,K,V)
=
\operatorname{softmax}
\left(
\frac{QK^\top}{\sqrt{d_k}}
\right)V.
\]

Hluboké modely se používají pro obraz, text, řeč, časové řady, fraud detection, doporučování nebo automatizaci
business procesů. Prakticky je důležité nejen trénování, ale také příprava dat, validace, MLOps, nasazení,
monitoring, škálování, bezpečnost a interpretovatelnost. Největší slabiny hlubokého učení jsou náročnost na
data a výpočetní výkon, riziko přeučení, black-box povaha a potřeba kontroly biasu, driftu a právních/etických
dopadů.

\newpage